\section{Well-founded relations and the Mostowski collapse}
\label{sec:cumulative-hierarchy-03}

\begin{definition}[Well-founded relation]
Let \(R\subseteq A\times A\).  The relation \(R\) is well-founded on \(A\)
iff every nonempty subset \(B\subseteq A\) has an \(R\)-minimal
member:
\[
  \forall B\subseteq A\,
  \bigl(B\ne\varnothing\to
        \exists b\in B\,\forall c\in B\,\neg(cRb)\bigr).
\]
It is extensional iff distinct elements have distinct sets of
\(R\)-predecessors.
\end{definition}

\begin{theorem}[Mostowski collapse]
Let \(A\) be a set and let \(R\subseteq A\times A\) be well-founded and
extensional.  There are a unique transitive set \(M\) and a unique isomorphism
\(\pi:(A,R)\to(M,\in)\).  They satisfy
\[
  \pi(a)=\set{\pi(b):bRa}
\]
for every \(a\in A\).
\end{theorem}

\paragraph{Proof sketch.}
Well-founded recursion---recursion along the well-founded set relation
\(R\), obtained from the set-recursion theorem via the rank function of
\(R\)---defines \(\pi\) by the displayed equation.  Replacement
makes each displayed collection a set.  An induction on \(R\) proves that
extensionality makes \(\pi\) injective.  Its range is transitive by the defining
equation, and the same equation proves preservation and reflection of the
relations.  A second well-founded induction proves uniqueness.  The
underlying recursion principle for arbitrary well-founded set relations and
the separation and replacement instances are supplied in
\cite[Chapter~6]{Jech2003}.

